\section{Legacy frozen-HLP source compatibility}
\label{sec:legacy-frozen-hlp}

This appendix records notation from the earlier frozen-HLP comparison because
the archived auxiliary Sections~\ref{sec:hlp-local} and
\ref{sec:threeZ} still refer to these labels.  None of the objects below
contributes to the principal operator or to the proof of
Theorem~\ref{thm:main}.

Put
\begin{equation}
\boxed{
\widehat H_{\rm ex}(s)
=f(s)-P(s)+\frac{Q(s)}{f(s)-P(s)}.
}
\label{eq:Hex-variable}
\end{equation}
On a smooth height block define
\begin{equation}
\boxed{F_b:=f(1/2+i\tau_b)\in\mathbb R,}
\label{eq:Fb-real-freeze}
\end{equation}
\begin{equation}
\boxed{H_{\rm ex}(s;F_b)=F_b-P(s)+\frac{Q(s)}{F_b-P(s)},}
\label{eq:Hex}
\end{equation}
\begin{equation}
\boxed{H_{\rm mod}(s;F_b)=F_b-2P(s)+P(s-2/L),}
\label{eq:Hmod}
\end{equation}
and
\begin{equation}
\boxed{
\Delta H(s;F_b)
=P(s)+\frac{Q(s)}{F_b-P(s)}-P(s-2/L).
}
\label{eq:deltaH}
\end{equation}
For compatibility with the old auxiliary $f'$ bookkeeping set
\begin{equation}
\boxed{q^2\left[\mathcal C\widehat H_{\rm ex}+R_{f'}^{\rm src}\right]}
\label{eq:exact-fprime-separation}
\end{equation}
with
\begin{equation}
\boxed{
R_{f'}^{\rm src}
=f'\zeta^2\widehat H_{\rm ex}
+\frac{f'}{f-P}\mathcal C
+\frac{(f')^2}{f-P}\zeta^2.
}
\label{eq:fprime-source}
\end{equation}
These are aliases for the superseded comparison only.

\begin{lemma}[Legacy locator for the translated model scale]
\label{lem:principal-pole-cluster}
Write $z=1+u$ and assume $F_b=L/2+O(1)$.  The frozen ratio
$H_{\rm ex}=F_b-P+Q/(F_b-P)$ has, as a meromorphic resummation near $u=0$,
a pole at $u=0$ of residue $-2$ and one zero-denominator pole
\[
u_*(F_b)=F_b^{-1}+O(F_b^{-2})=\frac2L+O(L^{-2})
\]
of residue $+1$.  The model $H_{\rm mod}$ has principal parts
$-2/u+1/(u-2/L)$.  This statement is not used by the revised theorem route.
\end{lemma}

\begin{proof}
Use
\[
P(1+u)=u^{-1}+p_0+O(u),
\qquad
Q(1+u)=u^{-2}+O(1).
\]
For $D(u)=F_b-P(1+u)$, Rouch\'e's theorem on a circle of radius
$O(F_b^{-2})$ about $F_b^{-1}$ gives one zero
$u_*=F_b^{-1}+O(F_b^{-2})$. Since $D'=Q$, the residue of $Q/D$ there is one.
At $u=0$, both $-P$ and $Q/(F_b-P)$ contribute residue $-1$.
\end{proof}
